\documentclass[a4paper,10pt,twocolumn]{ltjarticle}
\usepackage{kaitoushu}

\begin{document}

\problemrange{5章：三角関数 — 練習問題3・A (p.169)}

\mondai{1}{$\alpha$，$\beta$ がともに鈍角で，$\tan\alpha=-\dfrac{1}{2}$，$\cos\beta=-\dfrac{4}{5}$ のとき，次の値を求めよ．}
\kotae{
まず $\sin\alpha$，$\cos\alpha$，$\sin\beta$，$\tan\beta$ を求める．

$1+\tan^{2}\alpha=\dfrac{1}{\cos^{2}\alpha}$ より
$1+\dfrac{1}{4}=\dfrac{1}{\cos^{2}\alpha}$，$\cos^{2}\alpha=\dfrac{4}{5}$

$\alpha$ は鈍角だから $\cos\alpha<0$．

$\cos\alpha=-\dfrac{2}{\sqrt{5}}$，
$\sin\alpha=\tan\alpha\cos\alpha=\left(-\dfrac{1}{2}\right)\!\left(-\dfrac{2}{\sqrt{5}}\right)=\dfrac{1}{\sqrt{5}}$

$\beta$ も鈍角で $\cos\beta=-\dfrac{4}{5}$ だから $\sin\beta>0$．

$\sin\beta=\sqrt{1-\dfrac{16}{25}}=\dfrac{3}{5}$，
$\tan\beta=\dfrac{3/5}{-4/5}=-\dfrac{3}{4}$

\textbf{(1)} $\sin(\alpha+\beta)=\sin\alpha\cos\beta+\cos\alpha\sin\beta$

$=\dfrac{1}{\sqrt{5}}\cdot\left(-\dfrac{4}{5}\right)
+\left(-\dfrac{2}{\sqrt{5}}\right)\cdot\dfrac{3}{5}
=-\dfrac{4}{5\sqrt{5}}-\dfrac{6}{5\sqrt{5}}$

$=-\dfrac{10}{5\sqrt{5}}=-\dfrac{2}{\sqrt{5}}=-\dfrac{2\sqrt{5}}{5}$

\textbf{(2)} $\cos(\alpha+\beta)=\cos\alpha\cos\beta-\sin\alpha\sin\beta$

$=\left(-\dfrac{2}{\sqrt{5}}\right)\!\left(-\dfrac{4}{5}\right)
-\dfrac{1}{\sqrt{5}}\cdot\dfrac{3}{5}
=\dfrac{8}{5\sqrt{5}}-\dfrac{3}{5\sqrt{5}}$

$=\dfrac{5}{5\sqrt{5}}=\dfrac{1}{\sqrt{5}}=\dfrac{\sqrt{5}}{5}$

\textbf{(3)} $\tan(\alpha-\beta)=\dfrac{\tan\alpha-\tan\beta}{1+\tan\alpha\tan\beta}$

$=\dfrac{-\frac{1}{2}-\left(-\frac{3}{4}\right)}
{1+\left(-\frac{1}{2}\right)\!\left(-\frac{3}{4}\right)}
=\dfrac{\frac{1}{4}}{\frac{11}{8}}
=\dfrac{1}{4}\cdot\dfrac{8}{11}=\dfrac{2}{11}$
}

\mondai{2}{$\sin\alpha=-\dfrac{2\sqrt{2}}{3}$，$\pi<\alpha<\dfrac{3}{2}\pi$ のとき，$\sin2\alpha$，$\cos2\alpha$，$\sin\dfrac{\alpha}{2}$，$\cos\dfrac{\alpha}{2}$ の値を求めよ．}
\kotae{
$\pi<\alpha<\dfrac{3}{2}\pi$（第3象限）だから $\cos\alpha<0$．

$\cos^{2}\alpha=1-\dfrac{8}{9}=\dfrac{1}{9}$，$\cos\alpha=-\dfrac{1}{3}$

\medskip
2倍角の公式より

$\sin2\alpha=2\sin\alpha\cos\alpha
=2\left(-\dfrac{2\sqrt{2}}{3}\right)\!\left(-\dfrac{1}{3}\right)=\dfrac{4\sqrt{2}}{9}$

$\cos2\alpha=1-2\sin^{2}\alpha=1-2\cdot\dfrac{8}{9}=-\dfrac{7}{9}$

\medskip
半角の公式を使う前に $\dfrac{\alpha}{2}$ の範囲を調べる．

$\pi<\alpha<\dfrac{3}{2}\pi$ の各辺を2で割って
$\dfrac{\pi}{2}<\dfrac{\alpha}{2}<\dfrac{3}{4}\pi$

これは第2象限だから
$\sin\dfrac{\alpha}{2}>0$，$\cos\dfrac{\alpha}{2}<0$．

$\sin^{2}\dfrac{\alpha}{2}=\dfrac{1-\cos\alpha}{2}
=\dfrac{1+\frac{1}{3}}{2}=\dfrac{2}{3}$

$\therefore \sin\dfrac{\alpha}{2}=\sqrt{\dfrac{2}{3}}=\dfrac{\sqrt{6}}{3}$

$\cos^{2}\dfrac{\alpha}{2}=\dfrac{1+\cos\alpha}{2}
=\dfrac{1-\frac{1}{3}}{2}=\dfrac{1}{3}$

$\therefore \cos\dfrac{\alpha}{2}=-\dfrac{1}{\sqrt{3}}=-\dfrac{\sqrt{3}}{3}$
}

\mondai{3}{次の等式を証明せよ．}
\kotae{
\textbf{(1)} $\dfrac{\sin(\alpha+\beta)}{\sin(\alpha-\beta)}
=\dfrac{\tan\alpha+\tan\beta}{\tan\alpha-\tan\beta}$

加法定理で展開し，分子・分母を $\cos\alpha\cos\beta$ で割る．

左辺 $=\dfrac{\sin\alpha\cos\beta+\cos\alpha\sin\beta}
{\sin\alpha\cos\beta-\cos\alpha\sin\beta}$

$=\dfrac{\dfrac{\sin\alpha\cos\beta}{\cos\alpha\cos\beta}
+\dfrac{\cos\alpha\sin\beta}{\cos\alpha\cos\beta}}
{\dfrac{\sin\alpha\cos\beta}{\cos\alpha\cos\beta}
-\dfrac{\cos\alpha\sin\beta}{\cos\alpha\cos\beta}}
=\dfrac{\tan\alpha+\tan\beta}{\tan\alpha-\tan\beta}=$ 右辺 \qed

\medskip
\textbf{(2)} $\tan\!\left(\dfrac{\pi}{4}+x\right)\tan\!\left(\dfrac{\pi}{4}-x\right)=1$

$\tan\dfrac{\pi}{4}=1$ に注意して正接の加法定理を使う．

$\tan\!\left(\dfrac{\pi}{4}+x\right)=\dfrac{1+\tan x}{1-\tan x}$，\quad
$\tan\!\left(\dfrac{\pi}{4}-x\right)=\dfrac{1-\tan x}{1+\tan x}$

辺々掛けると

左辺 $=\dfrac{1+\tan x}{1-\tan x}\cdot\dfrac{1-\tan x}{1+\tan x}=1=$ 右辺 \qed
}

\mondai{4}{次の3倍角の公式を証明せよ．}
\kotae{
$3\theta=2\theta+\theta$ とみて加法定理と2倍角の公式を使う．

\textbf{(1)} $\sin3\theta=3\sin\theta-4\sin^{3}\theta$

$\sin3\theta=\sin(2\theta+\theta)=\sin2\theta\cos\theta+\cos2\theta\sin\theta$

$=2\sin\theta\cos\theta\cdot\cos\theta+(1-2\sin^{2}\theta)\sin\theta$

$=2\sin\theta\cos^{2}\theta+\sin\theta-2\sin^{3}\theta$

$\cos^{2}\theta=1-\sin^{2}\theta$ を代入して

$=2\sin\theta(1-\sin^{2}\theta)+\sin\theta-2\sin^{3}\theta$

$=2\sin\theta-2\sin^{3}\theta+\sin\theta-2\sin^{3}\theta$

$=3\sin\theta-4\sin^{3}\theta$ \qed

\medskip
\textbf{(2)} $\cos3\theta=4\cos^{3}\theta-3\cos\theta$

$\cos3\theta=\cos(2\theta+\theta)=\cos2\theta\cos\theta-\sin2\theta\sin\theta$

$=(2\cos^{2}\theta-1)\cos\theta-2\sin\theta\cos\theta\cdot\sin\theta$

$=2\cos^{3}\theta-\cos\theta-2\cos\theta\sin^{2}\theta$

$\sin^{2}\theta=1-\cos^{2}\theta$ を代入して

$=2\cos^{3}\theta-\cos\theta-2\cos\theta(1-\cos^{2}\theta)$

$=2\cos^{3}\theta-\cos\theta-2\cos\theta+2\cos^{3}\theta$

$=4\cos^{3}\theta-3\cos\theta$ \qed
}

\mondai{5}{積を和・差に直す公式を用いて，次の式を簡単な形にせよ．}
\kotae{
積和の公式

$\sin A\cos B=\dfrac{1}{2}\{\sin(A+B)+\sin(A-B)\}$

$\sin A\sin B=\dfrac{1}{2}\{\cos(A-B)-\cos(A+B)\}$

\textbf{(1)} $\sin\theta\cos3\theta+\sin\theta\cos5\theta+\sin\theta\cos7\theta$

各項を直すと（$\sin(-x)=-\sin x$ に注意）

$\sin\theta\cos3\theta=\dfrac{1}{2}(\sin4\theta-\sin2\theta)$

$\sin\theta\cos5\theta=\dfrac{1}{2}(\sin6\theta-\sin4\theta)$

$\sin\theta\cos7\theta=\dfrac{1}{2}(\sin8\theta-\sin6\theta)$

辺々加えると $\sin4\theta$ と $\sin6\theta$ が打ち消し合って

$=\dfrac{1}{2}(\sin8\theta-\sin2\theta)$

\textbf{(2)} $\sin\theta\sin3\theta+\sin\theta\sin5\theta+\sin\theta\sin7\theta$

同様に（$\cos(-x)=\cos x$ に注意）

$\sin\theta\sin3\theta=\dfrac{1}{2}(\cos2\theta-\cos4\theta)$

$\sin\theta\sin5\theta=\dfrac{1}{2}(\cos4\theta-\cos6\theta)$

$\sin\theta\sin7\theta=\dfrac{1}{2}(\cos6\theta-\cos8\theta)$

辺々加えて

$=\dfrac{1}{2}(\cos2\theta-\cos8\theta)$
}

\mondai{6}{次の式を1つの三角関数で表せ．}
\kotae{
$a\sin x+b\cos x=\sqrt{a^{2}+b^{2}}\,\sin(x+\varphi)$

ここで $\cos\varphi=\dfrac{a}{\sqrt{a^{2}+b^{2}}}$，
$\sin\varphi=\dfrac{b}{\sqrt{a^{2}+b^{2}}}$．

\textbf{(1)} $-\sqrt{3}\sin x+\cos x$

$a=-\sqrt{3}$，$b=1$ だから
$\sqrt{a^{2}+b^{2}}=\sqrt{3+1}=2$

$\cos\varphi=-\dfrac{\sqrt{3}}{2}$，$\sin\varphi=\dfrac{1}{2}$ より
$\varphi=\dfrac{5}{6}\pi$

$\therefore -\sqrt{3}\sin x+\cos x=2\sin\!\left(x+\dfrac{5}{6}\pi\right)$

\textbf{(2)} $3\cos x+\sqrt{3}\sin x=\sqrt{3}\sin x+3\cos x$

$a=\sqrt{3}$，$b=3$ だから
$\sqrt{a^{2}+b^{2}}=\sqrt{3+9}=2\sqrt{3}$

$\cos\varphi=\dfrac{\sqrt{3}}{2\sqrt{3}}=\dfrac{1}{2}$，
$\sin\varphi=\dfrac{3}{2\sqrt{3}}=\dfrac{\sqrt{3}}{2}$ より
$\varphi=\dfrac{\pi}{3}$

$\therefore 3\cos x+\sqrt{3}\sin x=2\sqrt{3}\sin\!\left(x+\dfrac{\pi}{3}\right)$
}

\mondai{7}{関数 $y=-\sin x+\cos x$ の最大値と最小値，およびそのときの $x$ の値を求めよ．ただし $0\leqq x<2\pi$ とする．}
\kotae{
合成する．$a=-1$，$b=1$ だから
$\sqrt{a^{2}+b^{2}}=\sqrt{2}$

$\cos\varphi=-\dfrac{1}{\sqrt{2}}$，$\sin\varphi=\dfrac{1}{\sqrt{2}}$ より
$\varphi=\dfrac{3}{4}\pi$

$y=\sqrt{2}\sin\!\left(x+\dfrac{3}{4}\pi\right)$

$0\leqq x<2\pi$ より
$\dfrac{3}{4}\pi\leqq x+\dfrac{3}{4}\pi<\dfrac{11}{4}\pi$

この範囲で $\sin$ が $1$，$-1$ になる角を探す．

\medskip
最大：$\sin\!\left(x+\dfrac{3}{4}\pi\right)=1$ となるのは
$x+\dfrac{3}{4}\pi=\dfrac{5}{2}\pi$（$\dfrac{\pi}{2}$ は範囲外）

$x=\dfrac{5}{2}\pi-\dfrac{3}{4}\pi=\dfrac{7}{4}\pi$

最小：$\sin\!\left(x+\dfrac{3}{4}\pi\right)=-1$ となるのは
$x+\dfrac{3}{4}\pi=\dfrac{3}{2}\pi$

$x=\dfrac{3}{2}\pi-\dfrac{3}{4}\pi=\dfrac{3}{4}\pi$

\medskip
$\therefore$ 最大値 $\sqrt{2}$（$x=\dfrac{7}{4}\pi$），
最小値 $-\sqrt{2}$（$x=\dfrac{3}{4}\pi$）
}

\end{document}
